
%%FILE blueprint/src/content.tex
% In this file you should put the actual content of the blueprint.
% It will be used both by the web and the print version.
% It should *not* include the \begin{document}
%
% If you want to split the blueprint content into several files then
% the current file can be a simple sequence of \input. Otherwise It
% can start with a \section or \chapter for instance.

\input{sections/topo}
\input{sections/unicite}
\input{sections/reduction}
\input{sections/rpf_theorem}
\input{sections/construction}


%%FILE blueprint/src/sections/construction.tex
\section{Construction d'une mesure de Gibbs}

  Soit $\lambda, \nu$ et $h$ comme dans le théorème 3.1.
  On pose $\mu = h \cdot \nu \in \Msp{n}$.

  \begin{lemma}
    \label{lem:hNu_sig_inv}
    \uses{def:shift, thm:rpf}
    \leanok
    \lean{RPF.mu_shift_inv, RPF.mu_prob}
    La mesure de probabilité $\mu = h \cdot \nu$ est invariante par $\sigma$ :
    $$\sigma_*\mu = \mu.$$
  \end{lemma}

  \begin{proof}
    \proves{lem:hNu_sig_inv}
    Soit $f, g \in \Csp{n}$, on remarque que $(\L f)g = \L (f\cdot(g\circ\sigma))$.
    En effet, pour $x \in \s_n$, on a
    \begin{align*}
      ((\L f) g)(x) &= \sum_{y\in\sigma^{-1}x}{e^{\phi(y)}f(y)g(x)} = \sum_{y\in\sigma^{-1}x}{e^{\phi(y)}f(y)g(\sigma y)} \\
		    &= \L(f\cdot(g\circ\sigma))(x).
    \end{align*}
    Donc pour $f \in \Csp{n}$,
    \begin{align*}
      \mu(f) &= \nu(h\cdot f) = \nu(\lambda^{-1}(\L h)\cdot f) = \lambda^{-1}\nu\left(\L(h\cdot(f\circ\sigma))\right) \\
	     &= \lambda^{-1}\L^*\nu(h\cdot(f\circ\sigma)) = \nu(h\cdot(f\circ\sigma)) = \mu(f\circ\sigma).
    \end{align*}
  \end{proof}

  A chaque fonction $f \in \Cs{n}$, on associe une nouvelle fonction $[f] \in \Csp{n}$ définie par :
  $$\forall x \in \sp_n, [f](x) = \min\left\{f(y)\mid y \in \s_n, \forall i \geq 0, x_i = y_i\right\}.$$
  De cette façon, on retire la dépendance de $f$ par rapport à ses coordonnées négatives.
  Par ailleurs, si $f \in \Csp{n}$, alors $[f] = f$.

  \begin{lemma}
    \label{lem:meas_seq_conv_to_sig_inv}
    \uses{lem:hNu_sig_inv} % + Riesz
    \leanok
    \lean{has_limit, has_limit_pos, G, G_one_eq_one}
    Pour tout $f \in \Cs{n}$, la suite $(\mu([f\circ\sigma^m]))_{m \geq 0}$ admet une limite, que l'on note $G(f)$.
    De plus $G$ vérifie $G(1) = 1$ et pour toute fonction $f \in \Cs{n}$, $f \geq 0$, on a $G(f) \geq 0$.

    Par ailleurs, $$G(f \circ \sigma) = G(f).$$
  \end{lemma}

  \begin{proof}
    \proves{lem:meas_seq_conv_to_sig_inv}
    \uses{lem:hNu_sig_inv}
    Soit $f \in \Cs{n}$, on va montrer que $(\mu([f\circ\sigma^n]))$ est une suite de Cauchy. Soit $m, k \in \N$,
    alors pour tout $x \in \sp_n$, il existe $y, y' \in \s_n$ tel que
    \begin{align*}
	[f\circ\sigma^m](\sigma^k x) &= f(\sigma^m(\cdots y_{-2}y_{-1}\textcolor{ForestGreen}{x_k}x_{k+1}\cdots))
		= f(\cdots y_{-1}x_k\cdots \textcolor{ForestGreen}{x_{m+k}}x_{m+k+1}\cdots), \\
	[f\circ\sigma^{m+k}](x) &= f(\sigma^{m+k}(\cdots y'_{-2}y'_{-1}\textcolor{ForestGreen}{x_0}x_1\cdots))
	    = f(\cdots y'_{-1}x_0\cdots\textcolor{ForestGreen}{x_{m+k}}x_{m+k+1}\cdots).
    \end{align*}
    Ainsi,
    $$\No{[f\circ\sigma^m]\circ\sigma^k - [f\circ\sigma^{m+k}]} \leq \var_k(f),$$
    et donc,
    $$\left| \mu([f\circ\sigma^m]) - \mu([f\circ\sigma^{m+k}]) \right| = \left|\mu\left([f\circ\sigma^m]\circ\sigma^k - [f\circ\sigma^{m+k}]\right) \right|
		\leq \var_k(f) \longrightarrow 0,$$
    lorsque $k \to \infty$, car $f$ est continue.
    On a alors prouvé que $(\mu([f\circ\sigma^m]))_m$ est une suite de Cauchy, ce qui assure sa convergence vers un réel $G(f)$.

    Comme $1 \in \Csp{n}$ et pour tout $m\in\N$ on a $1\circ\sigma^m = 1$, d'où
    $$\mu([1 \circ\sigma^m]) = \mu(1) = 1.$$

    Enfin, si $f \in \Cs{n}$ est positive, alors $[f] \geq 0$ et donc pour tout $m \in \N, [f \circ \sigma^m] \geq 0$,
    d'où $G(f) \geq 0$.
  \end{proof}

  Grâce au lemme précédent, on a construit une forme linéaire $G$ et par l'identification rendue possible par le théorème de Riesz,
  on peut considérer la mesure $\tilde{\mu} \in \Mss{n}$ associée à $G$.

  \begin{theorem}
    \label{thm:existence}
    \uses{lem:meas_seq_conv_to_sig_inv, def:gibbs_meas, def:holder_fn_set}
    \leanok
    \lean{is_gibbs}
    La mesure $\tilde\mu$ est une mesure de Gibbs pour $\phi \in \Hsp{n}$.
  \end{theorem}

  Pour prouver ce résultat, nous allons avoir besoin d'un lemme supplémentaire. Pour ce faire, on défini la constante $a$ donnée par :
  $$a = \sum_{k=0}^{\infty}{\var_k(\phi)} < \infty.$$

  \begin{lemma}
    \label{lem:birkhoff_sum_ineq}
    \uses{def:shift, def:topo_sn}
    \leanok
    \lean{birkhoff_ineq}
    Soit $x, y \in \s_n$ et $m \in \N$ tels que $y \in C(x, m)$.
    Alors,
    $$\left| S_m\phi(x) - S_m\phi(y)\right| \leq a.$$
  \end{lemma}

  \begin{proof}
    \proves{lem:birkhoff_sum_ineq}
    Pour $y \in \s_n$, on définit $y' \in \s_n$ par
    $$y'_i = \left\{\begin{array}{ll}
	x_i &\text{ si } i < 0, \\
	y_i &\text{ sinon.}
    \end{array}\right.$$
    Pour $k \geq 0$, on a alors $\phi(\sigma^ky) = \phi(\sigma^ky')$ car $\phi \in \Csp{n}$. Ainsi,
    \begin{align*}
      \left|S_m\phi(x) - S_m\phi(y)\right| = \left|S_m\phi(x) - S_m\phi(y')\right| &\leq \sum_{k=0}^{m-1}{\left|\phi(\sigma^kx) - \phi(\sigma^ky')\right|} \\
										   &\leq \sum_{k=0}^{m-1}{\var_{m-1-k}(\phi)} \leq a.
    \end{align*}
  \end{proof}

  \begin{proof}[Preuve du théorème \ref{thm:existence}]
    \proves{thm:existence}
    \uses{lem:birkhoff_sum_ineq, lem:meas_seq_conv_to_sig_inv}
    Soit $E = C(x,m) = \left\{y \in \s_n \mid \forall i \in \llbracket 0, m - 1\rrbracket, x_i = y_i\right\}$, et on veut montrer que
    $$c_1 e^{-Pm+S_m\phi(x)} \leq \tilde\mu(E) \leq c_2 e^{-Pm + S_m\phi(x)},$$
    pour certaines constantes $c_1, c_2 > 0$ et $P \in \R$.

    Pour tout $z \in \sp_n$, il existe un unique $y' = x_0x_1\cdots x_{m-1}z_0z_1\cdots \in \sp_n$ tel que $\sigma^my' = z$
    et $y'_i = x_i$ pour $0 \leq i \leq m - 1$. Ainsi,
    $$\L^m(h\ind{E})(z) = \sum_{\sigma^m y = z}{e^{S_m\phi(y)}h(y)\ind{E}(y)} \leq e^{S_m\phi(y')}h(y') \leq e^{S_m\phi(x)}e^a \No{h},$$
    et donc, comme $\ind{E} \in \Csp{n}$, on a $\tilde\mu(E) = \mu(E)$, d'où
    $$\tilde\mu(E) = \nu(h\ind{E}) = \lambda^{-m}\nu(\L^m(h\ind{E})) \leq \lambda^{-m}e^{S_m\phi(x)}\No{h}e^a.$$
    On peut alors poser $c_2 = e^a\No{h} > 0$.
    Pour l'autre inégalité, remarquons que pour $z \in \sp_n$, il existe au moins un $y' = x_0x_1\cdots 1 z_0z_1\cdots \in \sp_n$ tel que
    $\sigma^{m+1}y' = z$ et $x_i = y'_i$ pour $i \in \llbracket 0, m-1\rrbracket$, ce qui donne
    \begin{align*}
      \L^{m+1}(h\ind{E})(z) &= \sum_{\sigma^{m+1}y=z}{e^{S_m\phi(y) + \phi(\sigma^m y)}h(y)\ind{E}(y)} \\
			    &\geq e^{S_m\phi(y') - \No{\phi}}h(y') \\
			    &\geq (\inf h)e^{-\No{\phi} - a}e^{S_m\phi(x)},
    \end{align*}
    et donc
    $$\tilde\mu(E) = \nu(h\ind{E}) = \lambda^{-(m+1)}\nu(\L^{m+1}(h\ind{E}))
    \geq \lambda^{-m}e^{S_m\phi(x)}\underbrace{(\lambda(\inf h)e^{-\No{\phi}-a})}_{c_2}.$$
    On a les inégalités souhaitées avec $P = \log{\lambda}$, ce qui montre que $\tilde\mu$ est bien une mesure de Gibbs.
  \end{proof}

  \begin{proposition}
    \label{prop:mu_mixing}
    \uses{def:mixing, thm:existence}
    \leanok
    \lean{mu_mixing}
    La mesure $\tilde\mu$ est mélangeante par rapport à $\sigma$ : pour tout boréliens $E, F \in \mathcal B(\s_n)$ on a
    $$\tilde\mu(E \cap \sigma^{-m} F) \underset{m\to\infty}{\longrightarrow} \tilde\mu(E)\tilde\mu(F).$$
  \end{proposition}

  \begin{proof}
    \proves{prop:mu_mixing}
    \uses{lem:rpf_3_LambdaCr, def:topo_sn}
    Soit $x \in \s_n$, et $f \in \Csp{n}$. On prouve par récurrence que
    $$\L^mf(x) = \sum_{\sigma^my=x}{e^{S_m\phi(y)}f(y)}.$$
    Ainsi, si $g \in \Csp{n}$, on a
    $$((\L^mf)\cdot g)(x) = \sum_{\sigma^my = x}{e^{S_m\phi(y)}f(y)g(\sigma^my)} = \L^m(f\cdot(g\circ\sigma^m)).$$

    Pour prouver que $\tilde\mu$ est mélangeante, il suffit de le montrer pour $E, F$ dans la base de la topologie de $\s_n$.
    Soit alors $a, b \in \s_n$ et $E = \{y \in \s_n \mid y_i = a_i, r \leq i \leq s \}$ et $F = \{y \in \s_n \mid y_i = b_i, r \leq i \leq s \}$.
    Comme $\tilde\mu$ est invariante par $\sigma$, on peut supposer $r, u \geq 0$, de telle sorte que $\ind{E}, \ind{F} \in \Csp{n}$.
    \begin{align*}
      \tilde\mu(E\cap\sigma^{-m}F) &= \tilde\mu(\ind{E}\ind{\sigma^{-m}F}) = \tilde\mu(\ind{E}(\ind{F}\circ\sigma^m)) \\
				   &= \nu(h\ind{E}(\ind{F}\circ\sigma^m)) \\
				   &= \lambda^{-m}\L^{*m}\nu(h\ind{E}(\ind{F}\circ\sigma^m)) \\
				   &= \nu(\lambda^{-m}\L^m(h\ind{E})\ind{F}).
    \end{align*}
    De là, on obtient grâce au lemme \ref{lem:rpf_3_LambdaCr} et au fait que $\ind{E} \in \Cr{s}$
    \begin{align*}
      \left|\tilde\mu(E\cap\sigma^{-m}F)-\tilde\mu(E)\tilde\mu(F)\right| &= \left|\mu(E\cap\sigma^{-m}F) - \nu(h\ind{E})\nu(h\ind{F})\right| \\
									 &= \left|\nu\left[(\lambda^{-m}\L^m(h\ind{E})-\nu(h\ind{E})h)\ind{F}\right]\right| \\
									 &\leq \No{\lambda^{-m}\L^m(h\ind{E}) - \nu(h\ind{E})h}\nu(F) \\
									 &\leq A\nu(h\ind{E})\nu(F)\beta^{m-s} \longrightarrow 0,
    \end{align*}
    lorsque $m \to \infty$, ce qui permet de conclure.
  \end{proof}

  Finalement, on a construit une mesure de Gibbs mélangeante donc ergodique. En vertu de la proposition \ref{prop:unique},
  cette mesure est l'unique mesure de Gibbs pour cette fonction de potentiel $\phi$.


%%FILE blueprint/src/sections/intro.tex
\section{Introduction}

  \begin{definition}[Mesure de Gibbs]
    Soit $\mu$ une mesure de probabilité $\sigma$-invariante sur $\s_n$.
    On dit $\mu$ est une mesure de Gibbs pour un potentiel $\phi \colon \s_n \longrightarrow \R$ s'il existe $P \in \R$ et $c_1, c_2 > 0$ tels que
    pour tout $x \in \s_n$ et $m \in \N$ on ait
    $$c_1 \leq \frac{\mu\{y \in \s_n \mid \forall i \in \llbracket 0, m-1 \rrbracket, x_i = y_i\}}{\exp{(-Pm + \sum_{k=0}^{m-1}\phi(\sigma^kx))}} \leq c_2.$$
  \end{definition}

  L'objectif est de démontrer le théorème suivant, dû à R. Bowen.

  \begin{theorem}
    Soit $\phi$ une fonction de potentiel hölderienne. Alors il existe une unique mesure de Gibbs pour cette fonction $\phi$.
  \end{theorem}

  Pour ce faire, on se ramène au cas où la fonction de potentiel $\phi$ ne dépend plus des coordonnées négatives.
  Ensuite, on considère l'opérateur de transfert $\L$ défini par
  $$\forall f \in \Csp{n}, \forall x \in \sp_n, \hspace{0.5em} \L f(x) = \sum_{y \in \sigma^{-1}x}f(y)e^{\phi(y)},$$
  où $\sp_n$ est l'ensemble des suites à valeurs dans $\llbracket 1, n\rrbracket$ et indéxées sur $\N$.

  Le théorème suivant établit que cet opérateur admet une mesure propre et une fonction propre.

  \begin{theorem}[Ruelle-Perron-Frobenius]
    Soit $\phi$ un potentiel et $\L$ l'opérateur de transfert.
    Alors il existe $\lambda > 0, \nu \in \Msp{n}$ et $h \in \Csp{n}, h > 0$ tels que :
    \begin{enumerate}
      \item $\nu$ vérifie $\L^*\nu = \lambda\nu$,
      \item $h$ vérifie $\L h = \lambda h$ et $\nu(h) = 1$,
      \item et pour toute fonction $g \in \Csp{n}$,
	$\lim_{m\to\infty}\No{\frac 1 {\lambda^m} \L^mg - \nu(g)h} = 0$.
    \end{enumerate}
  \end{theorem}

  Pour prouver ce théorème, on utilisera le théorème de Schauder-Tychonoff afin de construire la mesure propre $\mu$
  et la fonction propre $h$ comme des points fixes de certains opérateurs, pour cela nous devrons d'abord établir la compacité de $\s_n$
  et d'un certain ensemble de fonctions notamment grâce au théorème d'Ascoli.
  Puis pour établir la limite nous aurons besoin de la densité des fonctions en escaliers dans $\Csp{n}$ et des propriétés de l'opérateur de transfert.

  Grâce à cette mesure propre $\nu$ et cette fonction propre $h$, on peut construire $\mu = h \cdot \nu$.
  Cette dernière mesure sur $\sp_n$ est alors $\sigma$-invariante, ce qui se montre grâce aux propriétés algébriques de l'opérateur de transfert et
  permettra de construire une forme linéaire $G$ sur $\s_n$ qui s'identifiera grâce au théorème de Riesz en une mesure $\tilde\mu$ sur $\s_n$,
  qui sera la mesure de Gibbs pour le potentiel höldérien $\phi$.
  Une fois $\tilde\mu$ construite, on montrera qu'elle est ergodique (et même mélangeante), ce qui permettra d'établir l'unicité.



%%FILE blueprint/src/sections/reduction.tex
\section{Réduction de $\s_n$ à $\sp_n$}

  Pour prouver l'existence de telles mesures, on commence par montrer que peu importe la fonction de potentiel $\phi$ dans $\Hs{n}$,
  on peut commencer par trouver une fonction de potentiel $\psi \in \Hsp{n}$
  ne dépendant que des coordonnées positives ayant la même mesure de Gibbs que $\phi$.

  \begin{definition}
    \label{def:fn_equiv}
    \uses{def:shift}
    \leanok
    \lean{Equiv.equiv}
    Soit $\phi, \psi \in \Cs{n}$. On dit que $\phi$ et $\psi$ sont équivalentes et on note $\phi \sim \psi$ dès lors
    qu'il existe $u \in \Cs{n}$ vérifiant :
    $$\phi = \psi - u + u \circ \sigma.$$
  \end{definition}

  Le prochain lemme justifie l'introduction de cette relation :

  \begin{lemma}
    \label{lem:fn_equiv_imp_gibbs_eq}
    \uses{def:fn_equiv, def:gibbs_meas}
    \leanok
    \lean{Equiv.equiv_imp_eq_gibbs}
    Soit $\phi \sim \psi \in \Hs{n}$. Dans ce cas, $\mu_{\phi} = \mu_{\psi}$ et $P(\phi) = P(\psi)$.
  \end{lemma}

  \begin{proof}
    \proves{lem:fn_equiv_imp_gibbs_eq}
    Soit $u \in \Cs{n}$ tel que $\phi - \psi = u\circ \sigma - u$. Alors
    \begin{align*}\left|S_m\phi(x) - S_m\psi(x)\right| &= \left|\sum_{k=0}^{m-1}(\phi - \psi)(\sigma^k x)\right|
							= \left|\sum_{k=0}^{m-1}(u(\sigma^{k+1}x) - u(\sigma^k x))\right| \\
						       &= \left|u(\sigma^m x) - u(x)\right| \leq 2 \lVert u \rVert
    \end{align*}

    Ainsi pour $m \in \N$ et $x \in \s_n$, on a
    $$c_1 e^{-2\lVert u\rVert} \leq \frac{\mu_{\phi}(C(x, m))}{e^{-P(\phi)m + S_m\phi(x) + 2\lVert u\rVert}}
			       \leq \frac{\mu_{\phi} (C(x, m))}{e^{-P(\phi)m + S_m\psi(x)}}
			       \leq \frac{\mu_{\phi}(C(x, m))}{e^{-P(\phi)m + S_m\phi(x) - 2\lVert u\rVert}} \leq c_2 e^{2\lVert u\rVert}$$

    Donc $\mu_{\phi}$ et $P(\phi)$ conviennent aussi pour $\psi$.
  \end{proof}

  \begin{remark}
    \label{rm:eq_if_equiv}
    Soit $\phi \overset{u}{\sim} \psi \in \Cs{n}$. Si $\sum_k{u\circ\sigma^k}$ converge,
    alors
    $$u = \sum_{k=0}^{\infty}{(\phi - \psi) \circ \sigma^k}$$

    Cette remarque permet de mieux comprendre la forme de la fonction $u$ dans un cadre favorable,
    et de donner l'intuition de ce qu'elle pourrait être pour donner une fonction équivalente à $\phi$ dépendant uniquement des coordonnées positives.
  \end{remark}

  \begin{proof}[Preuve de la remarque précédente]
    \proves{rm:eq_if_equiv}
    La convergence de $\sum_{k \geq 0}{(u - u\circ\sigma)\circ\sigma^k} = \sum_k{(\phi - \psi)\circ\sigma^k}$ est assuré
    par celle de $\sum_k{u\circ\sigma^k}$.
    Ainsi par téléscopage,
    $u(x) = \sum_{k=0}^{\infty}{(u - u \circ \sigma)(\sigma^k x)} = \sum_{k=0}^{\infty}{(\phi - \psi)(\sigma^k x)}$.
  \end{proof}

  Avant le prochain lemme on introduit la fonction $r \colon \s_n \longrightarrow \s_n$ définie par
  $$\forall i \in \Z, \forall x \in \s_n, (r(x))_i = \left\{
    \begin{array}{ll}
      x_i &\text{ si } i \geq 0, \\
      1   &\text{ si } i < 0.
    \end{array} \right. = 111\hdots x_{|\geq 0}$$
    De sorte que si $x, y \in \s_n$ ont les mêmes coordonnées positives (\textit{ie.} $\forall i \geq 0, x_i = y_i$), alors $r(x) = r(y)$.

    De plus, lorsqu'on compose $r$ avec $\sigma$, on garde l'indépendance vis-à-vis des coordonnées négatives :
    si $x = \cdots x_{-2}x_{-1}\textcolor{ForestGreen}{x_0} x_1 x_2 x_3\cdots \in \s_n$, alors on a :
    \begin{align*}
        x           = \cdots x_{-2}x_{-1}   &\textcolor{ForestGreen}{x_0} x_1 x_2\cdots \\
	\sigma x    = \cdots\cdot x_{-1}x_0 &\textcolor{ForestGreen}{x_1}x_2 x_3\cdots \\
	r(\sigma x) = \cdots\cdots\cdot 111 &\textcolor{ForestGreen}{x_1}x_2x_3\cdots \\
	r(x)	    = \cdots\cdots\cdot 111 &\textcolor{ForestGreen}{x_0}x_1 x_2\cdots \\
	\sigma r(x) = \cdots\cdots 11x_0    &\textcolor{ForestGreen}{x_1}x_2 x_3\cdots
    \end{align*}
    Donc $\sigma \circ r$ et $r \circ \sigma$ diffèrent seulement à la coordonnées -1 et ne dépendent que des coordonnées positives de $x$.

  \begin{lemma}
    \label{lem:fn_equiv_pos_coords}
    \uses{def:fn_equiv, def:holder_fn_set}
    \leanok
    \lean{Equiv.DependOnlyPosCoords, Equiv.equiv_pos_coords}
    Soit $\phi \in \Hs{n}$, et
    $$\psi := \phi \circ r + \left(\sum_{k=0}^{\infty}{\phi \circ \sigma^k}\right)\circ(\sigma \circ r - r \circ \sigma).$$
    Alors,
    \begin{enumerate}
      \item $\psi \in \Hs{n}$ et ne dépend que des coordonnées positives,
      \item Si on note $u = \sum_{k\geq 0}{\phi\circ\sigma^k - \phi\circ(\sigma^k r)}$, alors $u \in \Cs{n}$,
      \item De plus, $\phi \overset{u}{\sim} \psi$.
    \end{enumerate}
  \end{lemma}

  \begin{proof}
    \proves{lem:fn_equiv_pos_coords}
    Soit $b \in \R$ et $\alpha \in ]0, 1[$ tels que $var_k\phi \leq b\alpha^k$.
    Pour $k \geq 0$, $\sigma^k x$ et $\sigma^k r(x)$ coïncident de $-k$ à $+\infty$ donc
    $\left|\phi(\sigma^k x) - \phi(\sigma^k r(x))\right| \leq b\alpha^k$ et ainsi comme $|\alpha| < 1$, la série converge normalement
    pour tout $x \in \s_n$. Donc $u$ est bien définie et est continue.

    Ensuite vérifions que $\psi = \phi - u + u \circ \sigma$. Pour $N \in \N$,
    \begin{align*}
      &\phi(x) - \sum_{k=0}^{N}{\left(\phi(\sigma^k x)    - \phi(\sigma^k r(x))\right)}
	       + \sum_{k=0}^{N}{\left(\phi(\sigma^{k+1}x) - \phi(\sigma^k r(\sigma x))\right)} \\
      &= \phi(r(x)) - \sum_{k=0}^{N-1}{\left(\phi(\sigma^{k+1} x) - \phi(\sigma^{k+1} r(x))\right)}
		    + \sum_{k=0}^{N}{\left(\phi(\sigma^{k+1}x) 	  - \phi(\sigma^k r(\sigma x))\right)} \\
      &= \phi(r(x)) + \underbrace{\phi(\sigma^{N+1}x) - \phi(\sigma^N r(\sigma x))}_{\leq b\alpha^N}
		    + \sum_{k=0}^{N-1}{\left(\phi(\sigma^{k+1}r(x)) - \phi(\sigma^k r(\sigma x)) \right)} \\
      &\longrightarrow \psi(x) \\
      &\longrightarrow \phi(x) - u(x) + u(\sigma x)
    \end{align*}
    lorsque $N \to \infty$. Ainsi $\psi = \phi - u + u\circ\sigma$ \textit{ie.} $\phi \overset{u}{\sim} \psi$.

    Soit $x \in \s_n$ et $y \in C(x, m)$. Soit $k \in \N$, distinguons plusieurs cas :
    \begin{itemize}
      \item Comme $\sigma^kx$ et $\sigma^kr(x)$ coïncident jusqu'à la $k$-ième coordonnée,
	$$\left|\phi(\sigma^kx) - \phi(\sigma^kr(x)) - (\phi(\sigma^ky) - \phi(\sigma^kr(y)))\right|
	      \leq 2b\alpha^k$$
      \item Si $k \leq \lfloor\frac m 2\rfloor$, on peut faire mieux :
	\begin{align*}
	  &\left|\phi(\sigma^kx) - \phi(\sigma^ky) + \phi(\sigma^kr(x)) - \phi(\sigma^kr(y)) \right| \\
	  &\leq \underbrace{\left|\phi(\sigma^kx) - \phi(\sigma^ky)\right|}_{\leq var_{m-k}\phi \leq b\alpha^{m-k}}
	  + \underbrace{\left|\phi(\sigma^kr(x)) - \phi(\sigma^kr(y))\right|}_{\leq b\alpha^{m-k}} \\
	  &\leq 2b\alpha^{m-k}
	\end{align*}
	car $\sigma^kx$ et $\sigma^ky$ coïncident sur $[-m - k, m - k]$
    \end{itemize}
    Donc en sommant chacun des termes,
    \begin{align*}
      \left|u(x) - u(y)\right|
	    \leq 2b\left(\sum_{k=0}^{\lfloor\frac m 2\rfloor}{\alpha^{m - k}} + \sum_{k > \lfloor \frac m 2 \rfloor}{\alpha^k}\right)
	    &\leq 4b\sum_{k > \lfloor \frac m 2 \rfloor}{\alpha^k} = \frac{4b}{1-\alpha}\alpha^{\lfloor\frac m 2\rfloor} \\
	    &\leq \frac{4b\alpha}{1 - \alpha}(\sqrt{\alpha})^m
    \end{align*}
    Ainsi $u \in \Hs{n}$ et donc $\psi \in \Hs{n}$ également.
  \end{proof}

  \begin{remark}
    On remarque que la construction de $\psi$ dépend entièrement de la fonction $r$, ce qui signifie qu'il existe d'autres fonctions $\psi_r \in \Hs{n}$
    ne dépendant que des coordonnées positives et équivalentes à $\phi$.
  \end{remark}


%%FILE blueprint/src/sections/rpf_theorem.tex
\section{Théorème de Perron-Frobenius de Ruelle}

  \begin{definition}
    \label{def:transfert_operator}
    \uses{def:shift}
    \leanok
    \lean{RPF.transfert}
    Soit $\phi \in \Hsp{n} = \Csp{n} \cap \Hs{n}$. On appelle opérateur de transfert ou opérateur de Ruelle l'application
    $\L_{\phi} \colon \Csp{n} \longrightarrow \Csp{n} $ définie par :
    $$ \forall f \in \Csp{n}, \forall x \in \s_n, \L_{\phi}f(x) = \sum_{\sigma y = x}{e^{\phi(y)}f(y)}.$$
  \end{definition}

  \begin{remark}
    A noter que le passage aux fonctions définies sur $\Csp{n}$ rend $\sigma$ non injective, et donc $\L$ a toujours $n$ termes dans la somme
    qui le définit. C'est grâce à cette restriction que l'opérateur de transfert $\L$ devient intéressant.
  \end{remark}

  Pour la suite on fixe $\phi \in \Hsp{n}$ et des constantes $b > 0, \alpha \in ]0, 1[$ vérifiant $\var_k\phi \leq b\alpha^k$, pour tout $k \in \N$.

  \begin{theorem}[Théorème de Perron-Frobenius de Ruelle]
    \label{thm:rpf}
    \uses{def:transfert_operator, def:holder_fn_set, def:meas_set, prop:rpf_1, prop:rpf_2, prop:rpf_3}
    \leanok
    \lean{RPF.RPF1, RPF.RPF2, RPF.RPF3}
    Soit $\phi \in \Hsp{n}$ et $\L = \L_{\phi}$. Alors il existe $\lambda > 0, \nu \in \Msp{n}$ et $h \in \Csp{n}, h > 0$ tels que :
    \begin{enumerate}
      \item $\nu$ vérifie $\L^*\nu = \lambda\nu$,
      \item $h$ vérifie $\L h = \lambda h$ et $\nu(h) = 1$,
      \item et pour toute fonction $g \in \Csp{n}$, on a
	$$\No{\frac 1 {\lambda^m} \L^mg - \nu(g)h} \underset{m\to\infty}{\longrightarrow} 0.$$
    \end{enumerate}
  \end{theorem}

  Pour prouver ce théorème, nous allons avoir besoin de quelques lemmes et notamment du théorème suivant dû à Schauder et Tychonoff.

  \begin{theorem}[Théorème de point fixe de Schauder-Tychonoff]
    \label{thm:schauder}
    \leanok
    \lean{schauder_tychonoff}

    Soit $K$ un sous-ensemble compact convexe d'un espace vectoriel topologique localement convexe.
    Alors toute application $G \colon K \longrightarrow K$ continue admet un point fixe.
  \end{theorem}

  \begin{proposition}
    \label{prop:rpf_1}
    \uses{def:meas_set, def:transfert_operator}
    \leanok
    \lean{RPF.RPF1, RPF.RPF1_explicit}
    Il existe $\nu \in \Msp{n}$ tel que $\L^*\nu = \lambda\nu$.
  \end{proposition}

  \begin{proof}
    \proves{prop:rpf_1}
    \uses{def:meas_set, thm:schauder} % + Riesz
    Comme $\L1(x) = \sum_{\sigma y = x}{e^{\phi(y)}} > 0$ pour tout $x \in \s_n$, on a nécessairement $\L^*\mu(1) > 0$ pour tout $\mu \in \Msp{n}$.
    On peut alors poser $G(\mu) = \frac 1 {\L^*\mu(1)}\L^*\mu \in \Msp{n}$ pour toute mesure $\mu \in \Msp{n}$,
    ce qui donne une application $G\colon \Msp{n} \longrightarrow \Msp{n}$,
    (on peut se permettre $\Csp{n}^* = \Msp{n}$ grâce au théorème de Riesz).
    Ainsi définie, $G$ est continue de $\Msp{n}$ dans lui-même. Or $\Msp{n}$ est un compact convexe.
    Par le théorème de Schauder-Tychonoff, on en déduit que $G$ admet un point fixe $\nu \in \Msp{n}$.
    On pose alors $\lambda = \L^*\nu(1)$ et on a la relation voulue : $\L\nu = \lambda\nu$.
  \end{proof}

  Pour la suite on notera pour tout $m \in \N$, les constantes
  $$B_m = \exp{\left(\sum_{k\geq m+1}{2b\alpha^k}\right)} \hspace{0.5em}\text{ et }\hspace{0.5em} K = \lambda B_0 e^{\lVert\phi\rVert}.$$

  On cherche à construire $h$ comme un point fixe vérifiant (2), pour ce faire on considère l'ensemble $\Lambda \subseteq \Csp{n}$
  définie par :
  $$\forall f \in \Csp{n}, f\in \Lambda \iff
  \left\{\begin{array}{l}
    f \geq 0, \\
    \nu(f) = 1, \\
    \forall m \in \N, \forall x \in \s_n, x' \in C(x, m), f(x) \leq B_m f(x').
  \end{array}\right.$$
  On remarque que $1 \in \Lambda$, ce qui assure que $\Lambda \not= \emptyset$.

  \begin{lemma}
    \label{lem:Lambda_compact}
    \uses{def:transfert_operator}
    \leanok
    \lean{RPF.Lambda_IsCompact, RPF.Lf_ge_invK}
    L'ensemble $\Lambda$ est compact. De plus si $f \in \Csp{n}$, alors $\inf \frac 1 \lambda \L f \geq K^{-1}$.
  \end{lemma}

  \begin{proof}
    \proves{lem:Lambda_compact}
    % Uses : Ascoli
    Pour ce faire on va utiliser le théorème d'Ascoli. Dans un premier temps, on va montrer que $\Lambda(x) \subseteq [0, K]$.
    Soit $f \in \Lambda$ et $x \in \sp_n$. Remarquons que pour $z \in \sp_n$ et $x_0 \in \llbracket 1, n \rrbracket$ arbitraire, on a
    \begin{align*}
      \frac 1 \lambda \L f(x) = \sum_{\sigma y = x}{e^{\phi(y)}f(y)} &\geq \lambda^{-1}e^{-\lVert\phi\rVert}f(x_0x) \\
								     &\geq \lambda^{-1}e^{-\lVert\phi\rVert}B_0^{-1}f(z) = \frac 1 K f(z).
    \end{align*}
    Donc comme $1 = \nu(f) = \nu(\frac 1 \lambda \L f) \geq \frac 1 K f(z)$, on en déduit que $\lVert f\rVert \leq K$.
    De plus, comme $\nu(f) = 1$, il existe $z \in \sp_n$ tel que $f(z) \geq 1$. En appliquant l'inégalité précédente à un tel $z$,
    on obtient $\inf{\frac 1 \lambda \L f} \geq \frac 1 K$

    Soit $f \in \Lambda$ et $x, x' \in \sp_n$ tels que $x' \in C(x, m)$ pour un certain $m \in \N$,
    alors $f(x) \leq B_mf(x')$ et $f(x') \leq B_mf(x)$. Donc
    $$\left|f(x) - f(x')\right| \leq (B_m - 1)\lVert f\rVert \leq (B_m - 1)K \underset{m\to\infty}{\longrightarrow} 0.$$
    Ainsi comme la majoration est indépendante de $f$, $\Lambda$ est équicontinue.
  \end{proof}

  \begin{proposition}
    \label{prop:rpf_2}
    \uses{def:transfert_operator, prop:rpf_1, lem:Lambda_compact}
    \leanok
    \lean{RPF.RPF2, RPF.RPF2_explicit}
    Il existe $h \in \Lambda$ tel que $h > 0$ et vérifie (2) (\textit{ie.} $\nu(h) = 1$ et $\L h = \lambda h$).
  \end{proposition}

  \begin{proof}
    \proves{prop:rpf_2}
    \uses{lem:Lambda_compact, thm:schauder}
    On pose $F \colon \Csp{n} \longrightarrow \Csp{n}$ définie par
    $$\forall f \in \Csp{n}, F(f) = \frac 1 \lambda \L f.$$
    On vérifie que $\Lambda$ est stable par $F$. Soit $f \in \Lambda$,
    il est clair que $\frac 1 \lambda \L f$ est positif car $f$ l'est. Aussi,
    $\nu(\frac 1 \lambda \L f) = \frac 1 \lambda \L^*\nu(f) = \nu(f) = 1$.
    Enfin, pour $x, x' \in \sp_n$ tels que $x' \in C(x,m)$ pour un certain $m \in \N$, on a pour tout $j \in \llbracket 1, n\rrbracket$
    $$e^{\phi(jx)}f(jx) \leq e^{\phi(jx') + b\alpha^{m+1}}B_{m+1}f(jx') \leq B_me^{\phi(jx')}f(jx').$$
    En sommant tous ces termes, on obtient l'inégalité voulue : $\L f(x) \leq B_m\L f(x')$.
    Par ailleurs, $F$ est continue car $ne^{\lVert\phi\rVert}$-lipschitizienne : pour tout $f,g \in \Csp{n}$ et $x\in\sp_n$, on a
    $$\left|\L f(x) - \L g(x) \right| \leq \sum_{\sigma y = x}{e^{\phi(y)}\lVert f-g\rVert} \leq ne^{\lVert\phi\rVert}\lVert f-g\rVert,$$
    et donc $\lVert\L f - \L g\rVert \leq ne^{\lVert\phi\rVert}\lVert f-g\rVert$.

    En appliquant le théorème de Schauder-Tychonoff à $F$ sur $\Lambda$, on obtient un point fixe $h \in \Lambda$
    qui vérifie alors $\nu(h) = 1$ et $\L h = \lambda h$. De plus, par le lemme précédent,
    $0 < K^{-1} \leq \inf{\lambda^{-1}\L h} = \inf{h}$, d'où $h > 0$.
  \end{proof}

  %Comme $\log{B_m}$ converge vers 0 lorsque $m \to \infty$, cette suite est bornée par un réel $L > 0$.
  %Ainsi $\log{B_m^{-1}}, \log{B_{m+1}e^{b\alpha^m}}$ et $\log{B_m}$ sont dans $[-L', L']$ où $L' = L + b$.
  %Soit $\Delta = \left\{(x,x)\mid |x|\leq L'\right\}$.
  %La fonction $f\colon (x,y) \mapsto \frac{e^x - e^y}{x-y}$ est continue et strictement positive sur $[-L', L']^2\setminus\Delta$
  %et prolongeable par continuité sur $\Delta$ en posant $f(x,x) = e^x.$ pour tout $x \in [-L', L']$.
  %Finalement, $f$ est continue sur le compact $[-L',L']^2$ donc $f$ est bornée par $u_1, u_2 > 0$ et on a pour tout $x,y \in [-L', L']$ tels que $x<y$,
  %$$u_1(x-y)\leq e^x - e^y\leq u_2(x-y).$$

  \begin{lemma}
    \label{lem:Lf_decomp}
    \uses{def:transfert_operator, prop:rpf_2}
    \leanok
    \lean{RPF.Lf_decomp, RPF.decomp, RPF.Lf_decomp_explicit}
    Il existe $\eta \in ]0, 1[$ tel que pour toute fonction $f \in \Lambda$, il existe $f' \in \Lambda$ vérifiant
    $$\frac 1 \lambda \L f = \eta h + (1-\eta)f'.$$
    De plus, $\eta \leq \min{(\frac{u_1}{u_2} \frac{1-\alpha}{4\No{h}K}, \frac 1 {K\lVert h\rVert})}$.
  \end{lemma}

  \begin{proof}
    \proves{lem:Lf_decomp}
    \uses{lem:Lambda_compact, prop:rpf_2} % + Compacité et continuité
    Soit $\eta \in ]0, 1[$ comme ci-dessus. Soit $f \in \Lambda$ et on pose la fonction $g = \frac 1 \lambda \L f - \eta h$.
    Alors si $\eta \leq \min{(\frac{u_1}{u_2} \frac{1-\alpha}{4\No{h}K}, \frac 1 {K\lVert h\rVert})}$,
    on a $\frac 1 {1-\eta}g \in \Lambda$.
    En effet, l'intégrale de $g$ vaut $\nu(g) = \nu(\lambda^{-1}\L f) -\eta\nu(h) = 1 - \eta$, car $f$ et $h$ sont dans $\Lambda$.

    Aussi $g \geq 0$ dès lors que $\eta \lVert h\rVert \leq \frac 1 K$ car
    $$g = \lambda^{-1}\L f - \eta h \geq \inf{\lambda^{-1}\L f} - \eta\lVert h\rVert \geq 0,$$
    d'après le lemme \ref{lem:Lambda_compact}.

    Finalement, pour que $\frac 1 {1 - \eta}g \in \Lambda$ on doit avoir $g(x) \leq B_mg(x')$ pour $x \in \sp_n$ et $x'\in C(x,m)$ pour $m \geq 0$,
    ce qui équivaut à :
    $$\eta(B_mh(x') - h(x)) \leq B_m \lambda^{-1}\L f(x') - \lambda^{-1}\L f(x).$$
    Or, $\L f(x) \leq B_{m+1}e^{b\alpha^{m+1}}\L f(x')$ (comme vu dans la preuve du lemme \ref{lem:Lambda_compact}).
    Une condition suffisante est
    $$\eta(B_m - B_m^{-1})\lVert h\rVert \leq (B_m - B_{m+1}e^{b\alpha^{m+1}})K^{-1},$$
    car,
    $$\left\{
    \begin{array}{l}
      \eta(B_mh(x') - h(x)) \leq \eta(B_mh(x') - B_m^{-1}h(x')) \leq \eta(B_m-B_m^{-1})\lVert h\rVert, \vspace{1.5em}\\
      (B_m - B_{m+1}e^{b\alpha^{m+1}})K^{-1} \leq (B_m - B_{m+1}e^{b\alpha^{m+1}})\lambda^{-1}\L f(x') \leq \lambda^{-1}(B_m\L f(x') - \L f(x)).
    \end{array}\right.$$
    On rappelle que pour tout compact $K \subseteq \R^2$ et tout $x,y \in K$, on a
    $u_1(x-y) \leq e^x - e^y \leq u_2(x-y)$, ce qu'on applique ici avec le compact $K = [-L, L]^2$ de sorte que $\pm\log B_m, \log{B_m e^{b\alpha^m}}$
    soient dans $[-L, L]$.

    Ainsi il suffit de vérifier
    $$\eta u_2\lVert h\rVert(\log{B_m} - \log{B_m^{-1}}) \leq K^{-1}u_1(\log{B_m} - \log{B_{m+1}} - b\alpha^{m+1}),$$
    qui est équivalent à
    $$\eta\lVert h\rVert u_2 \left(\frac{4b\alpha^{m+1}}{1-\alpha}\right) \leq K^{-1}u_1b\alpha^{m+1},$$
    ou encore
    $$\eta \leq \frac{u_1}{u_2}\frac{1-\alpha}{4\lVert h\rVert K},$$
    afin d'avoir $\frac 1 {\eta - 1}g \in \Lambda$.
  \end{proof}

  Le lemme suivant est le cas particulier du théorème \ref{thm:rpf} pour les fonctions dans $\Lambda$, pour lesquelles en plus de la limite
  on a une convergence exponentielle.

  \begin{lemma}
    \label{lem:conv_expo}
    \uses{def:transfert_operator, prop:rpf_1, prop:rpf_2}
    \leanok
    \lean{RPF.conv_expo}
    Pour tout $n \in \N$ et tout $f\in\Lambda$, on a $\No{\lambda^{-n}\L^nf - h} \leq (\lVert h\rVert + K)(1 - \eta)^n = A\beta^n$
    avec $0 < \beta < 1$ et $A > 0$.
  \end{lemma}

  \begin{proof}
    \proves{lem:conv_expo}
    \uses{lem:Lf_decomp}
    Remarquons qu'en itérant le lemme précédent,
    \begin{align*}
      \lambda^{-1}\L f   &= \eta h + (1 - \eta)f'_1 = (1 - (1 - \eta)^1)h + (1 - \eta)f_1', \vspace{1em}\\
      \lambda^{-2}\L^2 f &= \lambda^{-1}\L((1 - (1 - \eta))h + (1 - \eta)f'_1) = (1 - (1 - \eta))h + \lambda^{-1}(1 - \eta)\L f'_1 \\
			 &= (1 - (1 - \eta) + \eta(1 - \eta))h + (1 - \eta)^2f'_2 \\
		         &= (1 - (1 - \eta)^2)h + (1 - \eta)^2f'_2,
    \end{align*}

    Ainsi, par récurrence et grâce au lemme précédent, on a pour tout $n \in \N$ et $f \in \Lambda$,
    $$\frac 1 {\lambda^n}\L^n f = (1 - (1 - \eta)^n)h + (1 - \eta)^nf_n,$$
    où $f_n \in \Lambda$. Comme $\lVert f_n\rVert \leq K$, on obtient
    $$\No{\frac 1 {\lambda^n}\L^n f - h} = (1-\eta)^n\lVert h + f_n\rVert \leq (1-\eta)^n(K + \lVert h\rVert).$$
  \end{proof}

  Avant d'étendre le résultat à toutes les fonctions de $\Csp{n}$, on l'étend d'abord à un sous-ensemble dense
  constitué des fonctions "en escaliers".
  On pose alors
  $$\Cr{r} = \left\{f \in \Csp{n} \mid \var_r(f) = 0 \right\} \text{ et } \mathcal C = \bigcup_{r\geq 0}{\Cr{r}}.$$
  On établira la densité de $\mathcal C$ dans un prochain lemme.

  \begin{lemma}
    \label{lem:Lambda_stb_mult_Cr}
    \uses{prop:rpf_1, lem:Lambda_compact}
    \leanok
    \lean{RPF.C_r, RPF.stab_mul_Cr}
    Soit $F \in \Lambda$ et $f \in \Cr{r}$ tels que $fF\not= 0$ et $f \geq 0$.
    Alors, $\frac{1}{\nu(fF)\lambda}\L^r(fF) \in \Lambda$.
  \end{lemma}

  \begin{proof}
    \proves{lem:Lambda_stb_mult_Cr}
    La positivité de $g = \lambda^{-r}\L^r(fF)$ découle de la positivité de $f$ et de $F$.
    Ensuite, soit $x, x' \in \sp_n$ et $m \in \N$ tels que $x' \in C(x,m)$.
    Remarquons d'abord que, grâce à une récurrence, on a
    $$\L^r(fF)(x) = \sum_{j_1, \hdots, j_r \in \llbracket 1, n\rrbracket}{\exp{\left(\sum_{k=0}^{r-1}{\phi(\sigma^k(j_1\hdots j_rx))}\right)}}
	  f(j_1\hdots j_r x)F(j_1\hdots j_r x).$$
    Soit $j_1, \hdots, j_r \in \llbracket 1, n\rrbracket$, alors comme $f \in \Cr{r}$ on a $f(j_1\hdots j_r x) = f(j_1\hdots j_r x')$. \\
    Aussi, $F(j_1\hdots j_r x) \leq B_{m+r} F(j_1\hdots j_r x')$ car les deux suites coïncident sur les $m+r$ premières coordonnées.
    Enfin,
    \begin{align*}
      B_{m+r}\exp{\left(\sum_{k=0}^{r-1}{\phi(\sigma^k(j_1\hdots j_r x))} \right)}
	&\leq B_{m+r} \exp{\left(\sum_{k=0}^{r-1}{\left(\var_{m+r-k}(\phi) + \phi(\sigma^k(j_1\hdots j_r x'))\right)}\right)} \\
	&\leq \left(B_{m+r} \exp{\left(\sum_{k = m+1}^{m+r} b\alpha^k\right)}\right) \exp{\left(\sum_{k=0}^{r-1}{\phi(\sigma^k(j_1\hdots j_r x'))}\right)} \\
	&\leq B_m \exp{\left(\sum_{k=0}^{r-1}{\phi(\sigma^k(j_1\hdots j_r x'))}\right)}.
    \end{align*}
    Ainsi, chacun des termes de $\L^r(fF)(x)$ est majoré par $B_m$ fois le terme correspondant dans $\L^r(fF)(x')$, d'où
    $\L^r(fF)(x) \leq B_m\L^r(fF)(x')$.

    Finalement, il reste à vérifier que $\nu(fF) > 0$. Pour ce faire, si $x, z \in \sp_n$ alors,
    $$\frac 1 \lambda \L(\L^r(fF))(x) = \lambda^{-1}\sum_{\sigma y=x}{e^{\phi(y)}\L^r(fF)(y)}
				      \geq \lambda^{-1}e^{-\No{\phi}}B_0^{-1}\L^r(fF)(z) = K^{-1}\L^r(fF)(z).$$
    Or comme $fF \not= 0$, il existe $z \in \sp_n$ tel que $(fF)(z) > 0$, donc $\L^r(fF)(\sigma^r z) > 0$. Ainsi,
    $$\nu(fF) = \frac 1 {\lambda^r}\nu(\lambda^{-1}\L(\L^r(fF))) \geq \frac 1 {K\lambda^r} \L^r(fF)(\sigma^r z) > 0.$$
    Et enfin, on a bien $\nu(\nu(fF)^{-1}\lambda^{-r}\L^r(fF)) = 1$, ce qui conclut ce lemme.
  \end{proof}

  \begin{lemma}
    \label{lem:rpf_3_LambdaCr}
    \uses{lem:conv_expo}
    \leanok
    \lean{RPF.RPF3_Lam_Cr}
    Soit $f \in \Cr{r}$, $F \in \Lambda$ et $n \in \N$.
    Alors,
    $$\No{\frac 1 {\lambda^{n+r}} \L^{n+r}(fF) - \nu(fF)h} \leq A\nu(|fF|)\beta^n.$$
  \end{lemma}

  \begin{remark}
    La fonction $F \in \Lambda$ n'a aucune utilité dans la preuve du théorème \ref{thm:rpf} (pour le montrer on prendra $F = 1$),
    elle sert plus tard pour montrer que $h\cdot\nu$ est la mesure de Gibbs du théorème \ref{thm:gibbs}.
  \end{remark}

  \begin{proof}
    \proves{lem:rpf_3_LambdaCr}
    \uses{lem:Lambda_stb_mult_Cr, lem:conv_expo}
    Décomposons $f$ en sa partie positive $f^+$ et sa partie négative $f^-$ toute deux positives et vérifiant $f = f^+ - f^-$.
    Si $f^{\pm}F \not= 0$, alors grâce aux lemmes \ref{lem:Lambda_stb_mult_Cr} et \ref{lem:conv_expo}, on a
    $$\No{\frac{1}{\lambda^{n+r}}\L^{n+r} (f^{\pm}F) - \nu(f^{\pm}F)h} = \nu(f^{\pm}F)\No{\frac{1}{\lambda^{n+r}\nu(f^{\pm}F)}\L^{n+r}(f^{\pm}F) - h}
	\leq A\nu(f^{\pm}F)\beta^n.$$
    Dans le cas où $f^{\pm}F = 0$ l'inégalité est triviale.
    Ainsi,
    $$\No{\frac 1 {\lambda^{n+r}} \L^{n+r} (fF) - \nu(fF)h} \leq A(\nu(f^+F) + \nu(f^-F))\beta^n = A\nu(|fF|)\beta^n.$$
  \end{proof}

  \begin{lemma}
    \label{lem:unionCr_strong_density}
    \uses{def:topo_sn}
    \leanok
    \lean{RPF.C_strong_density}
    L'ensemble $\mathcal C = \bigcup_{r\geq 0}{\Cr{r}}$ vérifie la propriété suivante :
    $$\forall f \in \Csp{n}, \forall \varepsilon > 0, \exists g_1, g_2 \in \mathcal C, \No{g_1 - g_2} \leq \varepsilon \text{ et } g_1\leq f\leq g_2,$$
    et donc $\mathcal C$ est dense dans $\Csp{n}$.
  \end{lemma}

  \begin{proof}
    \proves{lem:unionCr_strong_density}
    \uses{prop:sn_compact}
    Soit $f \in \Csp{n}$ et $\varepsilon > 0$. Alors il existe $r \in \N$ tel que pour tout $x \in \sp_n$ et $x' \in C(x, r)$ on ait
    $\left|f(x) - f(x')\right| \leq \varepsilon$, car $f$ est continue sur le compact $\sp_n$ donc uniformément continue par le théorème de Heine.
    On pose alors pour $x \in \sp_n$,
    $$\left\{\begin{array}{l}
	g_1(x) = \inf\limits_{z \in \sp_n}{f(x_1\cdots x_r z)}, \\
	g_2(x) = \sup\limits_{z \in \sp_n}{f(x_1\cdots x_r z)},
    \end{array}\right.$$
    de sorte que $g_1 \leq f \leq g_2$ et $g_1, g_2 \in \Cr{r} \subseteq \mathcal C$.
    De plus, si on note $z, z' \in \sp_n$ tel que $g_1(x) = f(x_1\cdots x_r z)$ et $g_2(x) = f(x_1 \cdots x_r z')$, alors
    $$\left|g_1(x) - g_2(x)\right| \leq \left|f(x_1\cdots x_r z) - f(x_1\cdots x_r z')\right| \leq \varepsilon.$$
  \end{proof}

  Enfin, on arrive au dernier lemme permettant d'établir le (3) du théorème \ref{thm:rpf}, pour lequel il ne reste plus que la limite
  sans la convergence exponentielle.

  \begin{proposition}
    \label{prop:rpf_3}
    \uses{def:transfert_operator, prop:rpf_1, prop:rpf_2}
    \leanok
    \lean{RPF.RPF3}
    Soit $f \in \Csp{n}$, alors
    $$\No{\frac 1 {\lambda^m} \L^mf - \nu(f)h} \underset{m \to \infty}{\longrightarrow} 0.$$
  \end{proposition}

  \begin{proof}
    \proves{prop:rpf_3}
    \uses{lem:unionCr_strong_density, lem:rpf_3_LambdaCr}
    Soit $\varepsilon > 0$, $r \in \N$ et $g_1, g_2 \in \Cr{r}$ comme dans le précédent lemme.
    On applique le lemme \ref{lem:rpf_3_LambdaCr} avec $F = 1$, ce qui donne pour $m$ assez grand,
    \begin{align*}
      \No{\frac 1 {\lambda^m}\L^m g_i - \nu(f)h} &\leq \No{\lambda^{-m}\L^mg_i - \nu(g_i)h} + \left|\nu(g_i) - \nu(f)\right|\No{h} \\
						 &\leq \varepsilon(1 + \No{h}).
    \end{align*}
    De plus, on a $\lambda^{-m}\L^mg_1 \leq \lambda^{-m}\L^mf \leq \lambda^{-m}\L^mg_2$ et donc pour $m$ assez grand,
    $$-\varepsilon(1 + \No{h}) \leq \lambda^{-m}\L^mg_1 - \nu(f)h
			       \leq \lambda^{-m}\L^mf - \nu(f)h \leq \lambda^{-m}\L^mg_2 - \nu(f)h
								\leq \varepsilon(1+\No{h}).
    $$
    Finalement, on a bien $\lim_{m\to\infty}{\No{\lambda^{-m}\L^mf - \nu(f)h}} = 0$, pour tout $f \in \Csp{n}$.
  \end{proof}

  \begin{remark}
    La densité de $\mathcal C$ dans $\Csp{n}$, à savoir
    $$\forall f \in\Csp{n}, \forall \varepsilon > 0, \exists r \in \N, \exists f_r\in \Cr{r}, \No{f - f_r} \leq \varepsilon,$$
    ne suffit pas pour conclure. En effet, si on a un telle fonction $f_r \in \Cr{r}$ pour un certain $r \in \N$,
    \begin{align*}
      \No{\frac 1 {\lambda^m}\L^mf  - \nu(f)h} &\leq \frac{1}{\lambda^m}\No{\L^mf - \L^mf_r}
						      + \No{\frac 1 {\lambda^m}\L^mf_r - \nu(f_r)h}
						      + |\nu(f_r) - \nu(f)|\No{h} \\
					    &\leq \frac{1}{\lambda^m}\No{\L^mf - \L^mf_r} + A\nu(|f_r|)\beta^{m - r} + \varepsilon\No{h}.
    \end{align*}
    Or, on n'arrive pas à majorer le premier terme, notamment car on ne possède aucune estimation de $\lambda$.
  \end{remark}


%%FILE blueprint/src/sections/topo.tex
\section{Topologie de $\s_n$}

  \subsection{Définition de l'espace métrique $\s_n$}

  \begin{definition}
    \label{def:topo_sn}
    \leanok
    \lean{Bernoulli}
    Soit $n \in \N$.
    L'ensemble des suites bi-infinies à valeurs dans $\llbracket 1, n\rrbracket$ est noté $\s_n = \prod_{\Z}{\llbracket 1, n\rrbracket}$.
    Chaque $\llbracket 1, n\rrbracket$ est muni de la topologie discrète,
    et on munit alors $\s_n$ de la distance produit $d$ donnée par
    $$\forall x, y \in \s_n,\hspace{1em} d(x, y) = \sum_{i\in\Z}{2^{-|i|}\ind{x_i \not= y_i}} < +\infty.$$
    On notera $B_d(x, r)$ la boule centrée en $x \in \s_n$ et de rayon $r \geq 0$ pour $d$.
  \end{definition}

  \begin{proposition}
    \label{prop:sn_compact}
    \uses{def:topo_sn}
    \leanok
    L'espace métrique $(\s_n, d)$ est compact.
  \end{proposition}

  \begin{proof}
    \proves{prop:sn_compact}
    \lean{Bernoulli}
    \leanok
    Chacun des $\llbracket 1, n\rrbracket$ est compact pour sa topologie.
    En tant que produit $\s_n$ est donc également compact, grâce au théorème de Tychonoff.
  \end{proof}

  \begin{remark}
    \label{rm:dist_equiv}
    Pour $x, y \in \s_n$ distincts, notons $N = \min\left\{i \in \N \mid x_i \not= y_i \text{ ou } x_{-i} \not= y_{-i}\right\}$.
    Alors $d_{\beta}(x, y) = \beta^N$ et $d_{\beta}(x, x) = 0$ définit une distance sur $\s_n$ et est équivalente à la distance $d$,
    et ce pour tout $\beta \in ]0, 1[$.
  \end{remark}

  \begin{proof}[Preuve de la remarque précédente]
    \proves{rm:dist_equiv}
    Soit $\beta \in ]0, 1[$. Alors $d_{\beta}$ est distance car $d_{\beta}$ est clairement symétrique, $d(x, y) = 0 \iff x = y$ pour tout $x,y\in\s_n$
    et $d_{\beta}$ vérifie l'inégalité triangulaire.

    De plus, $d_{\beta}$ est équivalente à $d$. En effet pour $r > 0$, montrons qu'il existe $r_1, r_2 > 0$ tel que pour tout $x \in \s_n$ on ait
    $$B_{\beta}(x, r_1) \subseteq  B_d(x, r) \subseteq B_{\beta}(x, r_2).$$
    Soit $y \in \s_n$. Supposons que $y \in B_d(x, r)$, alors $d(x, y) < r$, en particulier
    $$\forall i \in \Z, \hspace{1em}2^{-|i|}\ind{x_i \not= y_i} \leq r$$
    Soit $i \in \Z$, si $2^{-|i|} > r$, alors nécessairement $x_i = y_i$ et $|i| < -\frac{r}{\log{2}}$,
    alors en posant $r_2 = \beta^{-\frac{r}{\log{2}}}$ on a $y \in B_{\beta}(x, r_2)$, ce qui prouve une des deux inclusions.

    Supposons désormais que $x \in B_{\beta}(x, r)$, c'est-à-dire $d_{\beta}(x, y) < r$,
    donc $x$ et $y$ coïncident sur les $m_r = \frac{r}{\log{\beta}}$ premières coordonnées.
    Ainsi,
    $$d(x, y) = \sum_{i\in\Z}{2^{-|i|}\ind{x_i \not= y_i}} \leq 2\sum_{i \geq m_r + 1}{2^i} = 2^{1 - m_r} = r_1.$$
    et finalement, $y \in B_d(x, r_1)$. Donc les distances sont équivalentes et les topologies associées à ces distances sont les mêmes.
  \end{proof}

  On notera donc dans la suite, $B_{\beta}(x, r)$ les boules de centre $x \in \s_n$ et de rayon $r > 0$ pour la distance $d_{\beta}$,
  et $B(x, r) = B_{\frac 1 2}(x, r)$.
  On peut alors noter que si $y \in B_{\beta}(x, r)$, si on note $m_r = \frac{r}{\log{\beta}}$, alors $x$ et $y$ coïncident sur les coordonnées entre
  $-m_r$ et $m_r$ :
  $$\forall i \in \llbracket -m_r, m_r \rrbracket, \hspace{1em} x_i = y_i.$$

  \subsection{Propriétés topologiques de $\s_n$}

  On notera $\Cs{n}$ l'ensemble des fonctions continues $f \colon \s_n \longrightarrow \R$.

  \begin{definition}
    \uses{def:topo_sn}
    \label{def:var_potentiel}
    \leanok
    Soit $\phi \in \Cs{n}$. La fonction de variation d'ordre $k$ est donné par :
    $$\var_k(\phi) = \sup\left\{|\phi(x) - \phi(y)| : \forall i \in \llbracket -k, k\rrbracket, x_i = y_i \right\}.$$
  \end{definition}

  \begin{definition}
    \label{def:holder_fn_set}
    \uses{def:var_potentiel}
    \lean{HolderLike}
    \leanok
    On notera $\Hs{n}$ l'ensemble des fonctions $\phi \in \Cs{n}$ pour lesquels il existe $\alpha \in ]0, 1[$ et $b > 0$ vérifiant
    $$\forall k \in \N, \hspace{1em} \var_k(\phi) \leq b\alpha^k.$$
  \end{definition}

  \begin{proposition}
    \label{prop:holder_imp_continuous}
    \uses{def:holder_fn_set}
    \lean{holder_imp_continuous}
    \leanok
    Soit $\phi \in \Hs{n}$, alors $\phi$ est continue sur $\s_n$.
  \end{proposition}

  %Les résultats présentés dans la section Préliminaires s'appliquent donc pour $\s_n$. Cependant on peut caractériser de manière plus précise
  %la forme des ouverts de $\s_n$.

  \begin{proposition}
    \label{prop:separable}
    \uses{def:topo_sn}
    \leanok
    Pour tout $m \in \N$, soit $T_m = \left\{x \in \s_n \mid \forall|i|>m, x_i=1\right\}$ et $T = \bigcup_{m\geq 0}{T_m}$.
    L'ensemble $T$ est dense dans $\s_n$ et est dénombrable. Ainsi, $\s_n$ est séparable.
  \end{proposition}

  \begin{proof}
    \proves{prop:separable}
    \uses{def:topo_sn}
    \leanok
    Soit $O$ un ouvert non vide et $x \in O$, alors il existe $r = 2^{-m+1} > 0$ tel que $B(x, r) \subseteq O$.
    Si on prend $y \in \s_n$ tel que si $i \in \llbracket -m, m \rrbracket$ alors $y_i = x_i$ et sinon $y_i = 1$,
    alors $y \in T_m$ et $y \in B(x, r)$, donc $O \cap T \not= \emptyset$. Ainsi $T$ est dense dans $\s_n$.

    Pour tout $m \in \N$, l'ensemble $T_m$ contient exactement $n^{2m+1}$ élements, \textit{ie.} $T_m$ est fini.
    Donc par union dénombrable, $T$ est dénombrable.
  \end{proof}

  \begin{proposition}
    \label{prop:boules}
    \uses{def:topo_sn}
    \leanok
    \lean{Bernoulli.open_iff_disjoint_union_ball}
    Soit $O$ un ouvert de $\s_n$. Pour tout $x \in O$, on note $r_x > 0$ tel que $B(x, r_x) \subseteq O$.
    Alors il existe $\mathcal R \subseteq \s_n$ dénombrable tel que
    $$O = \bigsqcup_{x\in\mathcal R}{B(x, r_x)},$$
    où $\bigsqcup$ dénote l'union disjointe.
  \end{proposition}

  \begin{proof}
    \proves{prop:boules}
    \uses{prop:separable}
    \leanok
    Soit $O$ un ouvert de $\s_n$, alors $O = \bigcup_{x \in O}{B(x, r_x)}$.
    Cette union n'est pas disjointe et indicée sur un ensemble fini. On introduit alors la relation :
    $$\forall x,y \in O, x\sim y \iff B(x, r_x) \cap B(y, r_y) \not= \emptyset.$$
    Cette relation est clairement une relation d'équivalence. De plus on a pour $x, y \in O$,
    $$x \sim y \iff B(x, r_x) \subseteq B(y, r_y) \text{ ou } B(y, r_y) \subseteq B(x, r_x).$$
    En effet, supposons que $B(x, r_x) \cap B(y, r_y) \not= \emptyset$, alors soit $z$ dans l'intersection, on note $m_x = -\frac{r_x}{\log 2}$
    et $m_y = -\frac{r_y}{\log 2}$.
    Sans perte de généralité, on peut supposer que $m_y \leq m_x$. On a alors
    $$\left\{\begin{array}{l}
	\forall i \in \llbracket -m_x, m_x \rrbracket, x_i = z_i, \vspace{0.5em}\\
      \forall i \in \llbracket -m_y, m_y \rrbracket, y_i = z_i.
    \end{array}\right.$$
    Alors, pour tout $i \in \llbracket -m_y, m_y \rrbracket, x_i = y_i$ et donc $B(x, r_x) \subseteq B(y, r_y)$.
    Ainsi, il existe un système de représentant $\mathcal R \subseteq O$ pour la relation $\sim$ tel que
    $$O = \bigcup_{x\in O}{B(x, r_x)} = \bigsqcup_{x\in\mathcal R}{B(x, r_x)}.$$

    Il reste à montrer que $\mathcal R$ est au plus dénombrable.
    Pour tout $x \in \mathcal R$, il existe $h(x) \in T$ tel que $h(x) \in B(x, r_x)$.
    L'application $h \colon \mathcal R \longrightarrow T$ ainsi définie est donc injective et $T$ est dénombrable,
    ainsi $\mathcal R$ est dénombrable.
  \end{proof}

  \begin{definition}
    \label{def:shift}
    \uses{def:topo_sn}
    \leanok
    \lean{Bernoulli.shift}
    L'application de décalage (\textit{shift} en anglais) $\sigma$ sur $\s_n$ est donnée par
    $$(\sigma x)_i = x_{i+1}$$
    pour tout $x \in \s_n$ et $i \in \Z$.
    Cette application est alors un homéomorphisme de $\s_n$.
  \end{definition}

  \begin{definition}
    \label{def:cylindres}
    \uses{def:topo_sn}
    \leanok
    \lean{Bernoulli.cylinder}
    Soit $x\in \s_n$. On définit les cyclindres de centre $x$ et de rayon $m \in \N$ par
    $$C(x, m) = \left\{y \in \s_n \mid \forall i \in \llbracket 0, m - 1 \rrbracket, x_i = y_i\right\}.$$
    de cette manière on peut les relier aux boules pour la distance $d_{\beta}$ aux cyclindres :
    $$C(x,2m+1) = \sigma^{-m}(B_{\beta}(\sigma^mx, \beta^m)).$$
  \end{definition}


%%FILE blueprint/src/sections/unicite.tex
\section{Unicité de la mesure de Gibbs}

  Dans un premier temps, on introduit quelques notations.

  \begin{definition}
    \label{def:meas_set}
    \leanok
    On note $\Ms{n}$ l'ensemble des mesures de probabilité sur $\s_n$, et $\Mss{n}$ le sous-ensemble de $\Ms{n}$
    constitué des mesures $\sigma$-invariantes.
  \end{definition}

  \begin{definition}
    \label{def:gibbs_meas}
    \uses{def:meas_set, def:shift}
    \leanok
    \lean{IsGibbs}
    Soit $\phi \in \Cs{n}$ et $\mu \in \Ms{n}$.
    On dit que $\mu$ est une mesure de Gibbs pour un potentiel $\phi$ si $\mu \in \Mss{n}$ et s'il existe $P(\phi) \in \R$ et $c_1, c_2 > 0$ tels qu'on ait
    $$\frac{\mu(C(x, m))}{\exp(-Pm + \sum_{k=1}^{m-1}{\phi(\sigma^k(x))})} \in [c_1, c_2]$$
    pour tout $x \in \s_n$ et $m \in \N$.
  \end{definition}

  Dans cette section, on montre la partie concernant l'unicité du théorème suivant.

  \begin{theorem}
    \label{thm:gibbs}
    \uses{def:holder_fn_set, def:gibbs_meas, thm:existence, prop:mu_mixing, prop:unique, lem:fn_equiv_pos_coords, lem:fn_equiv_imp_gibbs_eq}
    \leanok
    \lean{gibbs, gibbs_unique}
    Soit $\phi \in \mathcal H(\s_n)$, alors il existe une unique mesure de Gibbs $\mu_{\phi} = \mu \in \mathcal M_{\sigma}(\s_n)$
  \end{theorem}

  \begin{proof}
    \proves{thm:gibbs}
    \uses{prop:unique, thm:existence}
    \leanok
    \lean{gibbs, gibbs_unique}
    La preuve découle des lemmes et propositions qui suivent.
  \end{proof}

  Dans les sections suivantes, on construira une mesure de Gibbs mélangeante notamment grâce au théorème de Ruelle.
  C'est l'ergodicité de cette mesure qui entrainera l'unicité de la mesure de Gibbs.

  \begin{definition}
    \label{def:ergodique}
    \uses{def:meas_set, def:shift}
    \leanok
    \lean{IsErgodic}
    Soit $\mu \in \Ms{n}$.
    On dit que $\mu$ est ergodique (par rapport à $\sigma$) si pour tout borélien $E \in \mathcal{B}(\s_n)$ tels que $\sigma^{-1}E=E$, on a
	$$\mu(E) = 0 \text{ ou } \mu(E^c) = 0,$$
  \end{definition}

  \begin{definition}
    \label{def:mixing}
    \uses{def:meas_set, def:shift}
    \leanok
    \lean{IsMixing}
    Soit $\mu \in \Ms{n}$. On dit que $\mu$ est mélangeante (par rapport à $\sigma$) si pour tout borélien $E, F \in \mathcal{B}(\s_n)$, on a
	$$\mu(E\cap\sigma^{-n}F) \underset{n\to\infty}{\longrightarrow} \mu(E)\mu(F).$$
  \end{definition}

  \begin{remark}
    On a l'implication "mélangeante" $\implies$ "ergodique".
  \end{remark}

  \begin{lemma}
    \label{lem:sig_inv_imp_cst}
    \uses{def:ergodique}
    \leanok
    \lean{ergodic_shift_inv_imp_cst}
    Soit $f \colon \s_n \longrightarrow \R$ intégrable par rapport à une mesure $\mu \in \Mss{n}$ ergodique.
    Supposons que $f \circ \sigma = f$ $\mu$-presque partout.
    Alors $f$ est constante $\mu$-presque partout.
  \end{lemma}

  \begin{proof}
    \proves{lem:sig_inv_imp_cst}
    Pour montrer ce lemme on considére les ensembles $E_c = f^{-1}(\{c\})$ pour tout $c\in\R$.
    Comme $f \circ \sigma = f$ $\mu$-p.p., on a $\sigma^{-1}E_c = E_c$, et donc par ergodicité de $\mu$,
    $$\mu(E_c) = 0 \text{ ou } \mu(E_c) = 1.$$
    Il est donc clair qu'il existe au plus un $c\in\R$ tel que $f = c$ pour $\mu$-presque tout $x \in \s_n$.
  \end{proof}

  \begin{proposition}
    \label{prop:unique}
    \uses{def:ergodique, def:gibbs_meas}
    \leanok
    \lean{Uniqueness.gibbs_ergodic_unique}
    Supposons qu'il existe une mesure de Gibbs $\mu$ ergodique pour $\phi \in \Cs{n}$.
    Alors cette mesure est l'unique mesure de Gibbs associé à $\phi$.
  \end{proposition}

  \begin{proof}
    \proves{prop:unique}
    \uses{lem:sig_inv_imp_cst, prop:boules} % + Radon-nikodym & régularité extérieure
    Soit $\mu, \mu' \in \Msp{n}$ deux mesures de Gibbs avec $\mu$ ergodique, $c_1, c_1', c_2, c_2' > 0$ et $P, P' \in \R$
    des constantes telles que pour tout $m \in \N$ et $x \in \s_n$ on ait
    $$\left\{\begin{array}{lll}
	c_1  &\leq \frac{\mu(C(x,m))}{\exp{(-Pm + S_m\phi(x))}} &\leq c_2 \vspace{1em} \\
	c_1' &\leq \frac{\mu'(C(x,m))}{\exp{(-P'm + S_m\phi(x))}} &\leq c_2'
      \end{array}\right.$$

      D'abord, montrons que $P = P'$.
      Soit $m \in \N$ et $T_m = \left\{x \in \s_n \mid x_i = 1 \text{ si } i \not\in \llbracket 0, m-1 \rrbracket \right\}$,
      de telle sorte que $\s_n = \bigsqcup_{x\in T_m}{C(x,m)}$, et $T_m$ est fini.
      On a alors
      $$c_1'e^{-P'm}\sum_{x\in T_m}{e^{S_m\phi(x)}} \leq \sum_{x\in T_m}{\mu'(C(x,m))}
            = 1 = \sum_{x\in T_m}{\mu'(C(x, m))} \leq c_2'e^{-P'm}\sum_{x\in T_m}{e^{S_m\phi(x)}},$$
      et donc
      $$P - \frac 1 m \log{c_2'} \leq \frac 1 m \log{\sum_{x\in T_m}{e^{S_m\phi(x)}}} \leq P - \frac 1 m \log{c_1'}.$$
      Par le théorème des gendarmes, on trouve donc $P' = \lim_{m\to\infty}{\frac 1 m\sum_{x\in T_m}{e^{S_m\phi(x)}}}$.
      Grâce au même raisonnement, on en conclut que $P = P'$ car ils sont tout deux égaux à la même limite.

      Grâce aux estimations sur $\mu$ et $\mu'$ sur les cylindres, on a alors pour tout $x \in \s_n$ et $m \in \N$,
      $$\mu'(C(x, m)) \leq \frac{c_2'}{c_1}\mu(C(x, m)).$$
      Comme $\mu$ et $\mu'$ sont invariantes par $\sigma$, on peut étendre le résultat sur les ensembles de la forme
      $\left\{y\in\s_n \mid \forall i \in \llbracket -m, m\rrbracket, x_i = y_i\right\}$ pour tout $x \in \s_n$ et $m \in \N$,
      c'est-à-dire sur une base de la topologie de $\s_n$.
      De plus, comme tout ouvert de $\s_n$ est une union disjointe dénombrable de ces éléments d'après la proposition \ref{prop:boules},
      on peut alors étendre cette inégalité aux ouverts de $\s_n$.
      Par régularité extérieure par rapport aux ouverts, on a pour tout borélien $B \in \Bo{\s_n}$,
      $$\mu'(B) = \inf\left\{\mu'(O) \mid O \text{ ouvert }, B \subseteq O\right\}.$$
      De là, on a pour tout ouvert $B \subseteq O$,
      $$\mu'(B) \leq \mu'(O) \leq \frac{c'_2}{c_1} \mu(O),$$
      et donc $\mu'(B) \leq \frac{c'_2}{c_1}\inf\{\mu(O) \mid O \text{ ouvert }, B \subseteq O\} = \frac{c_2'}{c_1}\mu(B)$.
      Donc pour tout borélien $B \in \Bo{\s_n}$ on a $\mu'(B) \leq \frac{c_2'}{c_1}\mu(B)$, ainsi $\mu'$ est absolument continue par rapport à $\mu$.
      D'après le théorème de Radon-Nikodym, $\mu'$ admet une densité $f$ par rapport à $\mu$.
      En appliquant $\sigma$, on obtient
      $$\mu' = \sigma_*\mu' = (f\circ\sigma^{-1})\cdot\sigma_*\mu = (f\circ\sigma^{-1})\cdot\mu.$$
      Par unicité de la dérivée de Radon-Nikodym, $f = f \circ \sigma^{-1}$ $\mu$-presque partout.
      Or comme $\mu$ est ergodique par hypothèse, alors il existe une constante $c \in \R$ tel que $f = c$ $\mu$-presque partout.
      Finalement,
      $$1 = \mu'(\s_n) = \int_{\s_n}{c\,d\mu} = c.$$
      Donc $\mu = \mu'$, ce qui prouve l'unicité.
  \end{proof}

